\chapter{The Axioms of a Mechanical Universe}
\label{ch:axioms}

% ======================================================================
%  PART I: THE ONTOLOGY
%  What exists, and what properties it has.
% ======================================================================

%----------------------------------------------------------------------
\section{Introduction: The Need for a New Foundation}
%----------------------------------------------------------------------

For over a century, theoretical physics has been defined by a profound and ever-widening schism.  At the grand scale, General Relativity describes a deterministic universe where gravity is the curvature of a spacetime manifold.  At the infinitesimal scale, the Standard Model describes a reality governed by probabilistic quantum fields and the exchange of virtual force-carriers.  These two pillars, despite immense predictive success, are fundamentally incompatible.

This schism has forced the acceptance of concepts that challenge mechanical intuition: superposition until observation, non-local action at a distance, and a vacuum teeming with virtual particles.  To reconcile theory with observation, we have postulated vast quantities of unseen dark matter and mysterious dark energy---placeholders for a deeper understanding we do not yet possess.

Spatial Displacement Theory (SDT) proposes a resolution by returning to first principles.  The universe is not abstract or probabilistic at its core; it is a deterministic, mechanical system governed by a clear and simple set of physical axioms.  SDT dismantles the need for curved spacetime, virtual particles, and acausal events by introducing a single, unified substrate for all of reality: a tangible, geometric, and pressurised spatial medium.

This chapter has two parts.  Part~I states the ontological axioms: \emph{what exists}.  Part~II derives the mathematical engine: \emph{how it computes}.


%----------------------------------------------------------------------
\section{The Foundational Constituents of Reality}
%----------------------------------------------------------------------

All complexity emerges from the interaction of three fundamental entities.

\subsection{Axiom I: Space (The Spation Medium)}

\textbf{Definition:} Space is not an empty void.  It is a tangible, physical medium---a discrete, densely packed, geometric lattice of zero-point, massless entities called \emph{spations}.

\textbf{Properties:} This medium is under immense, isotropic background pressure ($P_0$).  This pressure is a fundamental property of the universe.  The medium is the conduit for all interactions and its intrinsic properties define the universal speed limit, $c$.

\subsection{Axiom II: Matter (The Displacement Vortex)}

\textbf{Definition:} Particulate matter is not a foreign object existing \emph{in} space.  It is a localised, stable, dynamic vortex or standing-wave pattern \emph{of} space itself.  Where a particle exists, the spations of the medium are displaced.

\textbf{Properties:}
\begin{itemize}
  \item A particle's most fundamental property is the \textbf{geometry of its displacement vortex}.
  \item Its \textbf{mass} is a direct measure of the total energy of this displacement.
  \item Its \textbf{charge} is a property of the vortex's interaction with the surrounding pressure field.
  \item Its \textbf{spin} is the intrinsic chirality (handedness) of the vortex's rotation.
\end{itemize}

\subsection{Axiom III: Movement (The Kinetic Imperative)}

\textbf{Definition:} There is no true state of rest.  All energy is kinetic.  Every particle, every vortex, and the spations of the medium itself are in a state of eternal, perpetual motion and oscillation.

\textbf{Properties:} What is typically called ``potential energy'' is, in SDT, the stored kinetic energy of the pressurised spation medium being constrained by a displacement vortex.  A force is the transfer of momentum, and work is the result of reconfiguring these patterns of motion.


%----------------------------------------------------------------------
\section{The Emergent Nature of Time}
%----------------------------------------------------------------------

\subsection{Axiom IV: Time as a Measure of Change}

\textbf{Definition:} Time is not a fundamental, independent dimension.  Time is an \textbf{emergent property}; it is a local, relational measure of the sequence of events.

\textbf{The ``Now'':} The universe exists only in a singular, ever-present state of ``Now.''  The past is the memory of previous geometric configurations; the future is a projection of the likely next configurations.

\textbf{The ``Tick'' of a Clock:} The ``tick'' of any clock---pendulum or caesium transition---is a count of a regular, repeating sequence of spatial reconfigurations.  A clock does not measure the flow of an external ``time''; it counts its own internal, cyclical changes.  Apparent time dilation is not a warping of a time dimension, but a change in the local rate of physical oscillations due to changes in the local spation pressure field.


%----------------------------------------------------------------------
\section{The Core Mechanic: Pressure and Occlusion}
%----------------------------------------------------------------------

\subsection{Axiom V: Force as a Pressure Gradient}

\textbf{Definition:} Matter, by displacing the spation medium, creates a pressure field around itself.  This pressure is highest at the particle's boundary and diminishes with distance.  All forces result from other particles moving along these pressure gradients.

\textbf{Mathematical Principle:}
\begin{equation}\label{eq:force-gradient}
  \mathbf{F} \propto -\nabla P
\end{equation}
A particle's motion is always directed along the path of least resistance in the local pressure landscape.

\subsection{Axiom VI: The Principle of Occlusion}

\textbf{Definition:} A particle's displacement field blocks, or ``occludes,'' the pressure fields of all other particles in the universe.  This is a geometric ``shadowing'' effect.

\textbf{The Origin of ``Attraction'':} The apparent attractive force of gravity is a direct result of this occlusion.  Two bodies mutually shield each other from the isotropic background pressure ($P_0$).  This creates a low-pressure zone between them, into which the higher ambient pressure from the outside pushes them.  \textbf{Gravity is a ``push'' force, not a ``pull.''}


% ======================================================================
%  PART II: THE MATHEMATICAL ENGINE
%  Canonical Engine v4.1 — Audit-Hardened Core
%  Every symbol has ONE meaning and NEVER changes role.
% ======================================================================

\vspace{1cm}
\noindent\rule{\textwidth}{0.8pt}
\begin{center}
{\Large\bfseries Part II: The Mathematical Engine}\\[6pt]
{\itshape Canonical Engine v4.1 --- Audit-Hardened Core}
\end{center}
\noindent\rule{\textwidth}{0.8pt}
\vspace{0.5cm}


%----------------------------------------------------------------------
\section{Primitive Quantities}
\label{sec:primitives}
%----------------------------------------------------------------------

The engine operates on exactly four measured observables and four derived quantities.  No others are admitted.

\textbf{Measured observables:}
\begin{itemize}
  \item $c$ --- propagation speed constant
  \item $z$ --- measured wavelength shift fraction
  \item $r$ --- evaluation radius
  \item $v$ --- circular orbital velocity at $r$
\end{itemize}

\textbf{Derived quantities:}
\begin{itemize}
  \item $k$ --- dimensionless orbital parameter
  \item $v_{\text{esc}}$ --- escape velocity
  \item $\omega$ --- angular speed
  \item $v_{\text{rot}}$ --- rotation velocity
\end{itemize}

No gravitational constant.  No mass variable.  No $GM$.


%----------------------------------------------------------------------
\section{Definition Layer}
\label{sec:definitions}
%----------------------------------------------------------------------

These are definitions of measured quantities, not physical claims:
\begin{equation}\label{eq:k-def}
  k \equiv \frac{c}{v}
\end{equation}
\begin{equation}\label{eq:z-def}
  z \equiv \frac{\Delta\lambda}{\lambda}
\end{equation}

The symbol $k$ has \textbf{one meaning} throughout the corpus: the ratio of $c$ to the local orbital velocity.


%----------------------------------------------------------------------
\section{The SDT Bridge Law}
\label{sec:bridge-law}
%----------------------------------------------------------------------

The single structural relation proposed by SDT:
\begin{equation}\label{eq:bridge}
  z = \left(\frac{v}{c}\right)^2
\end{equation}

Therefore:
\begin{equation}\label{eq:z-from-k}
  z = \frac{1}{k^2}
\end{equation}

and:
\begin{equation}\label{eq:closure}
  \boxed{z \cdot k^2 = 1}
\end{equation}

This relation connects spectroscopy and orbital kinematics through one dimensionless closure condition.

\textbf{Epistemic note:} When $z$ is defined as $v^2/c^2$ and $k$ as $c/v$, the product $z \cdot k^2 = 1$ is an algebraic identity (see Chapter~2, \S2.3).  The \emph{physical} content of the bridge law is the claim that the \emph{spectroscopic observable} $\Delta\lambda/\lambda$ equals the kinematic quantity $v^2/c^2$.  This is a testable prediction, not a tautology.


%----------------------------------------------------------------------
\section{Orbital Consequences}
\label{sec:orbital}
%----------------------------------------------------------------------

From $k = c/v$:
\begin{equation}\label{eq:v-from-k}
  v = \frac{c}{k}
\end{equation}

Escape speed (standard virial relation):
\begin{equation}\label{eq:vesc}
  v_{\text{esc}} = \sqrt{2}\,v
\end{equation}

Angular velocity:
\begin{equation}\label{eq:omega}
  \omega = \frac{v}{r}
\end{equation}


%----------------------------------------------------------------------
\section{Radius Scaling}
\label{sec:radius-scaling}
%----------------------------------------------------------------------

Evaluating velocity at radius $r$ from a body of characteristic radius $R$:
\begin{equation}\label{eq:v-of-r}
  v(r) = \frac{c}{k}\,\sqrt{\frac{R}{r}}
\end{equation}

The parameter $k$ encodes the central body's displacement depth.  The $\sqrt{R/r}$ dependence reproduces Keplerian $v \propto r^{-1/2}$ scaling automatically.


%----------------------------------------------------------------------
\section{Spectral Formulation}
\label{sec:spectral}
%----------------------------------------------------------------------

Using the bridge law:
\begin{equation}\label{eq:v-from-z}
  v = c\sqrt{z}
\end{equation}

and:
\begin{equation}\label{eq:k-from-z}
  k = \frac{1}{\sqrt{z}}
\end{equation}

Thus spectroscopy alone determines the same parameter controlling orbital motion.  A measurement of wavelength shift is simultaneously a measurement of the displacement field depth.


%----------------------------------------------------------------------
\section{Rotational Coupling (Stellar Case)}
\label{sec:rotation}
%----------------------------------------------------------------------

Proposed SDT rotation relation for bodies in fusion equilibrium:
\begin{equation}\label{eq:rotation}
  v_{\text{rot}} = \frac{\pi\,v^2}{c}
\end{equation}

Angular rotation:
\begin{equation}\label{eq:omega-rot}
  \omega = \frac{v_{\text{rot}}}{r}
\end{equation}

This connects orbital motion to body rotation through a velocity-squared coupling.  The factor $\pi$ arises from the steradian geometry of the displacement field (Chapter~7).


%----------------------------------------------------------------------
\section{The Closure Hierarchy}
\label{sec:hierarchy}
%----------------------------------------------------------------------

Different physical systems correspond to different $k$ values:

\begin{center}
\begin{tabular}{lrl}
\toprule
\textbf{System} & $k$ & \textbf{Regime} \\
\midrule
Hydrogen ground state & 137 & Atomic \\
Sun surface & 686.5 & Stellar \\
Jupiter surface & 7\,124 & Planetary \\
Saturn surface & 11\,949 & Planetary \\
Earth surface & 37\,848 & Terrestrial \\
\bottomrule
\end{tabular}
\end{center}

The hierarchy follows from $z = 1/k^2$: deeper displacement levels (larger $k$) correspond to smaller spectral shifts.

The \textbf{atomic bridge constant} $\kop \equiv \alpha^{-1}\sqrt{R_p/a_0} = 0.5464$ (Chapter~3) sits at the bottom of this hierarchy, connecting nuclear geometry to atomic orbitals.  It is the $k$-value that bridges the proton charge radius to the Bohr radius through the fine structure constant.


%----------------------------------------------------------------------
\section{Interpretation}
\label{sec:interpretation}
%----------------------------------------------------------------------

In SDT:
\begin{itemize}
  \item $k$ represents the \textbf{dimensionless depth of a displacement field}.
  \item $z$ measures the \textbf{spectral signature of that depth}.
  \item Orbital motion follows directly from that depth through $v = c/k$.
\end{itemize}

Spectroscopy and orbital motion describe the same geometric quantity.  The bridge law $z = (v/c)^2$ is the single structural statement that unifies them.


%----------------------------------------------------------------------
\section{Comparison with Legacy Notation}
\label{sec:legacy}
%----------------------------------------------------------------------

Standard celestial mechanics expresses orbital strength through a fitted constant with units $\text{m}^3\text{s}^{-2}$ (the gravitational parameter $GM$).  In SDT, the equivalent observable quantity is:
\begin{equation}\label{eq:v2r}
  v^2 r
\end{equation}
which contains the same measurable information without introducing mass or a gravitational constant.  The equivalence is:
\begin{equation}\label{eq:gm-equiv}
  GM = v^2 r = \frac{c^2 R}{k^2}
\end{equation}

This is not a rejection of $GM$ as a \emph{computational tool}.  It is the observation that SDT's primitive variables ($c$, $v$, $r$, $k$) are sufficient, and $GM$ is a derived convenience, not a fundamental input.


% ======================================================================
%  CONCLUSION
% ======================================================================

%----------------------------------------------------------------------
\section{Conclusion}
%----------------------------------------------------------------------

This chapter presents the complete foundation of Spatial Displacement Theory in two layers:

\textbf{Part~I} (Axioms~I--VI) establishes the ontology: a pressurised spatial medium, displacement vortices, kinetic universality, emergent time, force as pressure gradient, and occlusion as the origin of apparent attraction.

\textbf{Part~II} (Canonical Engine v4.1) derives the mathematical structure from those axioms using exactly eight variables ($c$, $z$, $r$, $v$, $k$, $v_{\text{esc}}$, $\omega$, $v_{\text{rot}}$) and one structural relation: $z = (v/c)^2$.  Everything else follows algebraically.

From these foundations alone, the subsequent chapters derive the hierarchy of forces, the structure of the periodic table, and the observed dynamics of the cosmos.

\textbf{The real battlefield} is not notation or bookkeeping.  It is empirical discrimination:
\begin{enumerate}
  \item Does $z = (v/c)^2$ hold for gravitational redshift measurements?
  \item Does $v_{\text{rot}} = \pi v^2/c$ hold for stellar rotation?
  \item Does the $k$-hierarchy predict new systems without calibration?
\end{enumerate}

The cleaned engine above simply prevents the discussion from collapsing into arguments about imported constants or symbol drift.
